%%%%%%%%%%%%%%%%%%%%%---Chapter2 Formatting---%%%%%%%%%%%%%%%%%%%%%%%%%%%%%%%%%%
\documentclass[Thesis.tex]{subfiles}
\begin{document}
\newpage
\setcounter{chapter}{1} %%Chapter Numbering
\titleformat{\chapter}[display]
{\sffamily\fontsize{27}{27}\bfseries\filleft}{\thechapter}{0pt}{{#1}}
\thispagestyle{empty}
\chapter{Background and Literature Review}\label{chap2:background}

This chapter reviews the three bodies of work that this thesis draws on and,
in part, argues against. The first is the characterisation of noise in
pre-fault-tolerant quantum processors: what the dominant error mechanisms are,
how they are modelled as quantum channels, and how they are measured. The
second is quantum error mitigation (QEM), the family of post-processing
techniques that trade circuit repetitions for reduced bias in estimated
observables, and in particular the branch of that literature which replaces an
analytic extrapolation with a regression model learned from data. The third is
quantum machine learning (QML), which supplies the hybrid models benchmarked
here --- fidelity-kernel methods and data-reuploading variational circuits ---
together with a substantial body of negative and cautionary results about when
such models can be expected to help. A final strand, treated separately,
concerns decoding for quantum error correction (QEC), where machine-learned
decoders are a mature and well-benchmarked subfield and therefore a useful
methodological mirror for the mitigation work.

The chapter is organised so that the research gap identified in
Section~\ref{sec2:gap} follows from the review rather than being asserted on
top of it. Three recurring weaknesses in the learning-based mitigation
literature are traced through Sections~\ref{sec2:qem}--\ref{sec2:qml}:
inconsistent and often weak baselines, the almost universal absence of
cost normalisation against the circuit evaluations a method consumes, and the
near-total absence of tests outside the training distribution.

\section{Noise in pre-fault-tolerant quantum processors}\label{sec2:noise}

\subsection{Channels, Kraus operators and the observable-level picture}

The standard description of an open quantum system is the completely positive
trace-preserving (CPTP) map. For a density operator $\rho$ on a
$2^{n}$-dimensional Hilbert space, any physical noise process can be written
in Kraus form \cite{nielsen2012book}
\begin{equation}
  \mathcal{E}(\rho) \;=\; \sum_{k} K_{k}\,\rho\,K_{k}^{\dagger},
  \qquad \sum_{k} K_{k}^{\dagger}K_{k} \;=\; \mathbb{I}.
  \label{eq2:kraus}
\end{equation}
Three special cases dominate the modelling of superconducting and trapped-ion
hardware. Global depolarising noise of strength $p$,
\begin{equation}
  \mathcal{E}_{\mathrm{dep}}(\rho) \;=\; (1-p)\,\rho \;+\; p\,\frac{\mathbb{I}}{2^{n}},
  \label{eq2:depol}
\end{equation}
is the workhorse effective model because deep random circuits drive many
realistic error processes towards it; amplitude damping with rate
$\gamma = 1 - e^{-t/T_{1}}$ has Kraus operators
$K_{0} = \mathrm{diag}(1,\sqrt{1-\gamma})$ and
$K_{1} = \sqrt{\gamma}\,\lvert 0\rangle\langle 1\rvert$; and pure dephasing is
the Pauli channel with $K_{0}=\sqrt{1-p_{z}}\,\mathbb{I}$,
$K_{1}=\sqrt{p_{z}}\,Z$. Pauli channels in general,
\begin{equation}
  \mathcal{E}_{\mathrm{Pauli}}(\rho) \;=\; \sum_{a \in \{I,X,Y,Z\}^{n}} p_{a}\, P_{a}\,\rho\,P_{a}^{\dagger},
  \label{eq2:pauli}
\end{equation}
are important out of proportion to their physical realism, because they are
efficiently learnable, efficiently invertible in the Pauli basis, and can be
enforced on hardware by randomised compiling \cite{wallman2016noise}, which
twirls coherent errors into stochastic Pauli errors at no additional two-qubit
gate cost.

What matters for this thesis is not the full channel but its effect on a single
scalar. For an observable $O$ and an ideal $n$-qubit state
$\rho_{\mathrm{ideal}} = U\lvert 0\rangle\langle 0\rvert U^{\dagger}$ produced
by a target circuit $U$, the noisy device returns
\begin{equation}
  \langle O\rangle_{\mathrm{noisy}} \;=\; \mathrm{Tr}\!\left[\,O\,\widetilde{\mathcal{E}}(\rho_{\mathrm{ideal}})\,\right]
  \;=\; \langle O\rangle_{\mathrm{ideal}} \;+\; b(\lambda),
  \label{eq2:bias}
\end{equation}
where $\widetilde{\mathcal{E}}$ is the composition of the per-layer noise
channels interleaved with the ideal layers and $b(\lambda)$ is a bias that
depends on an effective noise strength $\lambda$. Under global depolarising
noise and a traceless observable, Eq.~\eqref{eq2:bias} collapses to the
particularly simple multiplicative form
$\langle O\rangle_{\mathrm{noisy}} = (1-p)\,\langle O\rangle_{\mathrm{ideal}}$.
Every mitigation method reviewed below is, at bottom, an estimator of
$\langle O\rangle_{\mathrm{ideal}}$ built from measurements at one or more
values of $\lambda$, and every one of them pays for reduced bias with increased
variance at fixed shot budget. The reviews by Preskill
\cite{preskill2018nisq}, Bharti et al.\ \cite{bharti2022noisy} and Cai et al.\
\cite{cai2023review} frame the same trade-off at the level of whole
algorithms.

\subsection{Characterisation: from averaged fidelities to learned noise models}

Mitigation quality is bounded by how well the noise is known, so the
characterisation literature is a direct input. Randomised benchmarking (RB)
estimates an average gate error from the exponential decay of survival
probability over random Clifford sequences of increasing length
\cite{knill2008randomized}; its appeal is that the extracted decay rate is
insensitive to state-preparation and measurement (SPAM) errors and requires no
tomographic reconstruction. RB in its standard form returns a single averaged
number and is therefore blind to the structure of the error --- it cannot
distinguish coherent over-rotation from stochastic Pauli noise of the same
average fidelity, a distinction that matters enormously for whether
extrapolation-based mitigation behaves smoothly in $\lambda$.

Gate set tomography (GST) addresses this by self-consistently estimating gates,
preparations and measurements up to gauge freedom, yielding full process
matrices rather than scalar fidelities \cite{nielsen2021gst}. GST is
informative but scales poorly, which is why the most consequential recent
development for mitigation is the efficient learning of \emph{structured}
sparse noise models. Harper, Flammia and Wallman showed that the Pauli error
rates of a full $n$-qubit Pauli channel can be estimated efficiently from
cycle-benchmarking-style experiments \cite{harper2020efficient}, and van den
Berg et al.\ turned this into a practical device-level pipeline by fitting a
sparse Pauli--Lindblad generator whose support is restricted to the
connectivity graph \cite{vandenberg2023probabilistic}. This learned-model route
is what makes probabilistic error cancellation deployable at the scale of a
hundred qubits, and it is the reason mitigation on current hardware is best
understood as a two-stage procedure: learn a noise model, then invert or
extrapolate through it. Learning-based mitigation of the kind benchmarked in
this thesis is a third option --- learn the map from noisy to ideal
expectation values directly, bypassing the explicit channel model --- and its
central empirical question is whether that shortcut generalises.

A practical corollary of the characterisation literature is that the effective
noise seen by a mitigation routine is partly a matter of compilation choice
rather than device physics. Coherent errors accumulate quadratically in
amplitude and are not described by a Pauli channel, so their contribution to
the bias $b(\lambda)$ in Eq.~\eqref{eq2:bias} is neither monotone nor smooth
in the number of repetitions of a gate; twirling them into stochastic Pauli
noise by randomised compiling restores the smooth, approximately exponential
$\lambda$-dependence on which both extrapolation and learned regression
implicitly rely \cite{wallman2016noise}. Two consequences follow for the
benchmark reported in this thesis. First, the noise model used for a simulation
study must be stated precisely, because the ranking of mitigation methods can
depend on whether the dominant error is coherent or stochastic. Second, an
empirical claim that a learned mitigator beats an analytic extrapolant is a
claim about a specific noise structure unless the experiment varies that
structure deliberately. The reviews of the NISQ programme by Preskill
\cite{preskill2018nisq} and Bharti et al.\ \cite{bharti2022noisy}, and of
mitigation specifically by Cai et al.\ \cite{cai2023review}, all treat the
noise model as the central modelling assumption rather than an implementation
detail; much of the learned-mitigation literature does not.

\section{Quantum error mitigation}\label{sec2:qem}

\subsection{Zero-noise extrapolation: estimator, bias and variance}

Zero-noise extrapolation (ZNE) was introduced independently by Temme, Bravyi
and Gambetta \cite{temme2017error} and, in the guise of active error
minimisation for variational simulators, by Li and Benjamin
\cite{li2017efficient}. The construction assumes that the expectation value
can be measured not only at the native noise strength $\lambda_{0}$ but also at
amplified strengths $\lambda_{i} = c_{i}\lambda_{0}$ with known scale factors
$c_{i} \ge 1$, and that the resulting family of estimates is a smooth function
of $\lambda$ admitting a series expansion
\begin{equation}
  \langle O\rangle(\lambda) \;=\; \langle O\rangle_{\mathrm{ideal}}
  \;+\; \sum_{k=1}^{m} a_{k}\,\lambda^{k} \;+\; \mathcal{O}(\lambda^{m+1}).
  \label{eq2:zne-series}
\end{equation}
Given $m+1$ measurements at distinct scale factors, the Richardson
extrapolant is the unique degree-$m$ polynomial interpolant evaluated at
$\lambda = 0$, which can be written as a fixed linear combination of the
measured values,
\begin{equation}
  \widehat{\langle O\rangle}_{\mathrm{ZNE}} \;=\; \sum_{i=0}^{m} \gamma_{i}\,\langle O\rangle(c_{i}\lambda_{0}),
  \qquad
  \gamma_{i} \;=\; \prod_{j \ne i} \frac{c_{j}}{c_{j}-c_{i}},
  \qquad \sum_{i} \gamma_{i} = 1 .
  \label{eq2:richardson}
\end{equation}
Two properties of Eq.~\eqref{eq2:richardson} govern everything that follows.
First, the residual bias is
$\mathcal{O}(\lambda_{0}^{m+1})$ only if the expansion~\eqref{eq2:zne-series}
is valid over the whole sampled range; when the true dependence is not
polynomial --- and for deep circuits it is typically closer to an exponential
decay --- Richardson extrapolation of increasing order does not converge to
the noiseless value but instead amplifies model misspecification. Second, the
coefficients $\gamma_{i}$ alternate in sign and grow with $m$, so the variance
of the estimator is inflated by
$\sum_{i}\gamma_{i}^{2}$ relative to a single measurement. Higher-order ZNE
therefore buys bias reduction at a quantifiable and rapidly worsening
statistical price, which is the reason practical implementations rarely go
beyond $m = 2$ or $3$.

Recognising the exponential character of noise accumulation, Endo, Benjamin and
Li fitted decaying exponentials rather than polynomials
\cite{endo2018practical}, and Cai generalised this to multi-exponential
ans\"atze with several decay constants, showing on both simulation and hardware
that a sum of exponentials tracks the observed $\lambda$-dependence of deep
circuits better than either a polynomial or a single exponential
\cite{cai2021multiexponential}. The comparison between linear, Richardson and
exponential extrapolation is thus a comparison of prior assumptions about the
functional form of $b(\lambda)$, evaluated at the same measurement cost; this
thesis retains all three as analytic baselines for exactly that reason.

\subsection{How noise is amplified}\label{sec:ch2_noise_amplification}

ZNE requires a knob for $\lambda$, and the three established knobs are not
equivalent. Kandala et al.\ implemented the original proposal by stretching
the microwave pulses that implement each gate, so that a gate of nominal
duration $t$ is executed over $ct$ while the rotation angle is held fixed
\cite{kandala2019extends}. This is the most faithful realisation of
``the same circuit under more noise'' because it scales the physical
decoherence exposure directly, and it was used to extend the reach of a
variational chemistry calculation on a superconducting device. It also demands
pulse-level control and an assumption that the noise per unit time is
stretch-invariant.

Digital alternatives replace pulse control with circuit rewriting.
Unitary folding maps each gate or each circuit layer $G \mapsto G G^{\dagger}
G$, which is logically the identity but triples the noise exposure;
Giurgica-Tiron et al.\ formalised local and global folding, showed how
fractional scale factors are obtained by folding a random subset of gates, and
benchmarked the resulting estimators \cite{giurgicatiron2020digital}. He et
al.\ studied the closely related strategy of inserting identity pairs
(\textit{e.g.}\ pairs of CNOTs) and pointed out that the effective noise
scaling achieved is only approximately linear in the number of insertions,
because inserted gates need not carry the same error profile as the gates they
are meant to emulate \cite{he2020zero}. The \textsc{Mitiq} package made these
scaling and extrapolation choices reproducible and interoperable across
back-ends \cite{larose2022mitiq}, which is why folding-based ZNE has become the
de facto baseline in the mitigation literature. The third knob, parameter-noise
scaling, perturbs gate rotation angles to emulate stronger error and is
likewise available as a standard scaling method \cite{larose2022mitiq}; it is
cheap but conflates coherent and incoherent contributions and is therefore
poorly suited to devices whose dominant error is stochastic.

The practical consequence, and one that the benchmarking design of this thesis
relies on, is that the ZNE estimator consumes a number of circuit evaluations
equal to the number of scale factors. A three-point extrapolation is three
circuit executions per estimated observable, not one. Any comparison between
ZNE and a learned mitigator that ignores this factor is comparing methods at
different budgets.

\subsection{Probabilistic error cancellation and noise-model inversion}

Probabilistic error cancellation (PEC) takes the opposite approach: rather than
extrapolating in a scalar, it inverts the learned noise channel exactly. If the
noise $\mathcal{E}$ is known, its inverse can be decomposed as a
quasiprobability combination of implementable operations
$\{\mathcal{O}_{i}\}$,
\begin{equation}
  \mathcal{E}^{-1} \;=\; \sum_{i} q_{i}\,\mathcal{O}_{i},
  \qquad q_{i} \in \mathbb{R},
  \qquad \sum_{i} q_{i} = 1,
  \qquad \gamma \;\equiv\; \sum_{i} \lvert q_{i}\rvert \;\ge\; 1,
  \label{eq2:pec}
\end{equation}
and an unbiased estimate of the ideal expectation value is obtained by sampling
$\mathcal{O}_{i}$ with probability $\lvert q_{i}\rvert/\gamma$ and reweighting
each shot by $\gamma\,\mathrm{sgn}(q_{i})$
\cite{temme2017error,endo2018practical}. PEC is unbiased when the noise model
is exact, and its cost is explicit: the variance is inflated by $\gamma^{2}$,
and $\gamma$ grows exponentially in the circuit volume. Van den Berg et al.\
made PEC practical at scale by fitting sparse Pauli--Lindblad models and
sampling the inverse of the learned generator
\cite{vandenberg2023probabilistic}, and Kim et al.\ combined learned-noise
inversion with amplification to obtain expectation values on a 127-qubit
processor for circuits beyond exact classical simulation
\cite{kim2023utility}. That experiment is the clearest statement of the
current state of the art and of its price: accuracy is bought with sampling
overhead that grows with the volume of the circuit being mitigated.

\subsection{Readout, symmetry and purification-based methods}

Measurement errors are a separately structured and comparatively benign error
source. If the assignment matrix $A$ with entries
$A_{ij} = \Pr(\text{observe }i \mid \text{prepared }j)$ is estimated from
calibration circuits, corrected probability vectors follow from
$\hat{p} = A^{-1} p_{\mathrm{meas}}$, usually with a constrained least-squares
solve to keep $\hat{p}$ a distribution. The obstacle is that $A$ has
$4^{n}$ entries; Nation et al.\ removed this obstacle with a
matrix-free formulation that restricts attention to the subspace of observed
bit strings and solves the system iteratively, making readout mitigation
routine at the scale of tens of qubits \cite{nation2021scalable}.

Symmetry verification exploits the fact that many circuits of interest conserve
a known quantity, such as particle number in a fermionic simulation.
Bonet-Monroig et al.\ showed that measuring the symmetry operator and
discarding or post-selecting runs that violate it removes a substantial
fraction of the error at very low cost, and that it composes with
extrapolation \cite{bonetmonroig2018symmetry}. A related family of methods
attacks the problem at the level of the state rather than the observable:
virtual distillation estimates
$\mathrm{Tr}[O\rho^{M}]/\mathrm{Tr}[\rho^{M}]$ from $M$ copies of the noisy
state, exponentially suppressing the contribution of non-dominant eigenvectors
of $\rho$ and thereby the incoherent error, at the cost of additional qubits
and entangling operations between copies \cite{huggins2021virtual}. These
methods are complementary to ZNE rather than competitive with it, and their
relevance here is that they illustrate the general pattern: each mitigation
technique defines its own resource currency --- extra scale factors, extra
shots with sign weights, extra copies, discarded runs --- and comparisons are
meaningful only once that currency is made explicit.

\subsection{Sampling cost and fundamental limits}

The theoretical literature has made the pessimistic side of this trade-off
precise. Takagi et al.\ derived general lower bounds on the sampling cost of
any error mitigation protocol that is required to produce an unbiased
estimator, showing that the overhead must grow exponentially in circuit depth
for broad classes of noise \cite{takagi2022fundamental}. Quek et al.\
subsequently tightened the picture considerably: for local depolarising noise
at constant rate, the number of circuit executions any mitigation strategy
needs grows exponentially with depth, and beyond a depth threshold the noisy
device becomes effectively useless for estimating observables no matter how
much classical post-processing is applied \cite{quek2024exponentially}. The
results are asymptotic statements about worst-case families and do not forbid
useful mitigation at the modest depths accessible today --- Kim et al.'s
experiment is the counterpoint \cite{kim2023utility} --- but they do settle the
question of whether mitigation can substitute for error correction: it cannot.
They also carry a methodological lesson for benchmarking. If the cost of
mitigation is the central theoretical quantity, then an empirical comparison
that reports only accuracy, with no accounting of the circuit evaluations
consumed, is measuring the less interesting half of the problem.

\subsection{Composition and the common currency of circuit evaluations}

The methods above are frequently composed: symmetry post-selection before
extrapolation \cite{bonetmonroig2018symmetry}, readout correction before
either \cite{nation2021scalable}, learned Pauli-channel inversion combined
with amplification \cite{kim2023utility}, or exponential fits layered on
purified estimates \cite{huggins2021virtual}. Composition is where the
accounting problem becomes acute, because each stage multiplies the sampling
overhead: if stage $j$ inflates the variance by a factor $\kappa_{j}$, the
total shot requirement to reach a target standard error scales as
$\prod_{j}\kappa_{j}$, and the reported accuracy of a composite pipeline is
meaningless without that product. Cai et al.\ organise the whole field around
exactly this bias--variance--overhead structure \cite{cai2023review}, and the
\textsc{Mitiq} implementation makes the composition explicit in software by
separating the noise-scaling executor from the inference step
\cite{larose2022mitiq}.

For the purposes of this thesis, the useful common currency is the number of
distinct circuit executions required per estimated expectation value, holding
the per-circuit shot count fixed. On this measure, the unmitigated estimate
costs one; an $m+1$-point extrapolation costs $m+1$ regardless of the
functional form fitted \cite{endo2018practical,cai2021multiexponential}; a
regressor whose input features are noisy expectation values at several noise
levels costs one per feature; and a hybrid quantum model costs the same as its
classical counterpart plus the executions needed to evaluate its feature map
or ansatz. Expressing every method in this currency is what makes the nine-way
comparison in this thesis interpretable, and it is what most of the studies
reviewed in the next section omit.

\section{Learning-based and machine-learning-based error mitigation}\label{sec2:mlqem}

The methods of Section~\ref{sec2:qem} share a common structure: a fixed
estimator, derived from an assumed functional form or an explicit noise model,
applied to measurements of the circuit of interest. Learning-based mitigation
replaces the fixed estimator with a function fitted to data. Write
$\mathbf{x}$ for a vector of features measured on the noisy device --- most
often the noisy expectation value itself, optionally together with values at
amplified noise, circuit metadata, or related observables --- and
$y = \langle O\rangle_{\mathrm{ideal}}$ for the target. The mitigator is a
regression function $f_{\theta}$ trained by empirical risk minimisation over
training pairs generated from circuits whose exact answers are known,
\begin{equation}
  \theta^{\star} \;=\; \arg\min_{\theta}\;
  \frac{1}{N}\sum_{i=1}^{N}\bigl(f_{\theta}(\mathbf{x}_{i}) - y_{i}\bigr)^{2}
  \;+\; \Omega(\theta),
  \label{eq2:erm}
\end{equation}
with $\Omega$ a regulariser. The whole approach therefore rests on two
questions that are statistical rather than quantum-mechanical: where do the
training labels $y_{i}$ come from, and over what distribution of circuits is
the fitted $f_{\theta^{\star}}$ expected to remain valid?

\subsection{Classically simulable training data: Clifford data regression and its descendants}

The answer to the first question that unlocked the field is near-Clifford
circuit construction. Czarnik et al.\ proposed Clifford data regression (CDR),
which builds training circuits that are structurally similar to the target
circuit but contain few enough non-Clifford gates to be simulated classically
by stabiliser methods, giving exact labels $X^{\mathrm{exact}}_{i}$ alongside
device measurements $X^{\mathrm{noisy}}_{i}$; a linear ansatz
$X^{\mathrm{exact}} \approx a X^{\mathrm{noisy}} + b$ is then fitted and
applied to the target circuit. They reported an order-of-magnitude error
reduction for a ground-state energy problem on 16 qubits of an IBMQ device and
on a 64-qubit noisy simulator, and analysed performance as a function of qubit
count, depth, and the number of non-Clifford gates \cite{czarnik2021clifford}.
Strikis et al.\ used the same near-Clifford substitution trick for a different
purpose: instead of fitting a scalar correction, they learn the
quasiprobability compensation strategy itself, so that PEC-style cancellation
becomes possible without a prior characterisation of the noise
\cite{strikis2021learning}. Lowe et al.\ then unified the two branches of
data-driven mitigation. Their variable-noise CDR (vnCDR) generates training
data both by near-Clifford substitution and by varying the noise level of those
training circuits, so that the fitted model has access to the same
$\lambda$-dependence that ZNE exploits. On an 8-qubit Ising energy estimation
task with a noise model taken from hardware they reported improvement factors
of 33 over the unmitigated result, 20 over ZNE and 1.8 over plain CDR, with
further benchmarks on 64-qubit random circuits \cite{lowe2021unified}.

Two features of this line of work deserve emphasis. First, the baselines are
explicit and the comparisons are like-for-like: vnCDR is compared against both
of the methods it unifies, on the same circuits and the same noise model.
Second, the training distribution is deliberately matched to the target ---
training circuits are built to resemble the circuit that will be mitigated ---
which means the method is designed around, rather than tested against, the
problem of distribution shift. A separate contribution in the same family,
the resource-efficient error suppression technique of Czarnik et al., uses
near-Clifford circuits not to fit a correction but to select the optimal number
of virtual copies in a purification scheme, reducing the qubit requirement from
$MN+1$ to $2N+1$ by exploiting active qubit reset
\cite{czarnik2021qubit}; it is a useful reminder that ``learning'' in this
literature sometimes means hyperparameter selection rather than function
approximation.

\subsection{Neural network regressors for expectation values}

The step from linear ans\"atze to flexible nonlinear regressors was taken
early. Kim, Park and Rhee proposed using an artificial neural network as the
mitigation map \cite{kim2020ann}, an early explicit statement of mitigation as
a supervised learning problem; it is cited here for priority. Zlokapa and
Gheorghiu applied learning one level lower in the stack, training a
convolutional network on device data to predict circuit noise and then using it
inside a compiler that inserts gate sequences minimising predicted noise; on
IBM 5-qubit devices this reduced output noise by $12.3\%$ relative to the
default compiler, and --- importantly for the generalisation question --- a
model trained on one device transferred to another with only $5.2\%$
improvement, demonstrating that the learned object was device-specific rather
than a universal noise map \cite{zlokapa2020deep}.

Bennewitz et al.\ introduced neural error mitigation for ground-state problems,
using neural quantum states to post-process variational quantum eigensolver
output for the H$_{2}$ and LiH molecular Hamiltonians and for the lattice
Schwinger model, and reported improved energies, fidelities and derived
quantities such as order parameters and entanglement entropy without
additional quantum resources \cite{bennewitz2022neural}. Zhukov and Pogosov
studied a deep network applied at the post-processing stage for Trotterised
dynamics of a two-dimensional spin lattice in the strong-noise regime, and
their negative result is as informative as their positive one: weight-one and
weight-two observables were recovered well under depolarising and
inhomogeneous Pauli noise, but coherent $ZZ$ crosstalk was \emph{not}
mitigated, so that randomised compiling \cite{wallman2016noise} is needed
first to convert coherent errors into stochastic ones
\cite{zhukov2024quantum}. This is direct evidence that a learned mitigator
inherits assumptions about noise structure even when no noise model is written
down.

Two proposals attack the training-data bottleneck itself. Babukhin et al.\
generate labelled pairs by echo evolution --- evolving a state forward and then
backward in time so that the ideal final state is known by construction, while
the device output carries the accumulated noise --- which removes the need for
either classical simulation or target-circuit simplification
\cite{babukhin2024echo}. Liao et al.\ go further and build a neural mitigator
that requires neither prior knowledge of the noise nor any noise-free training
data, using a quantum augmentation construction, and demonstrate it on
simulated circuits and on hardware \cite{liao2025noiseagnostic}. In the
measurement-error setting, Lee et al.\ exploited conditional independence of
distant qubits to factorise the correction and combined it with transfer
learning, obtaining an exponential reduction in network size relative to a
monolithic deep model and validating on IBM devices
\cite{lee2023scalable} --- the machine-learning analogue of the structural
simplification that made matrix-free readout mitigation scalable
\cite{nation2021scalable}. Sack and Egger pushed learned mitigation into an
application setting, using a feed-forward network to mitigate \emph{samples}
rather than expectation values so that depth-two QAOA could be optimised on
non-planar random regular graphs with up to 40 nodes, requiring circuits with
958 two-qubit gates \cite{sack2024large}.

\subsection{The closest prior work, and what it measured}

The most directly comparable study to this thesis is that of Liao et al.\
\cite{liao2024machine}, who benchmarked a family of machine-learning
mitigators --- linear regression, random forests, multi-layer perceptrons and
graph neural networks --- on hardware using up to 100 qubits, across several
circuit classes and increasingly complex device-noise profiles, with digital
zero-noise extrapolation as the reference method, and who explicitly examined
both interpolation and extrapolation regimes. Their headline claim is a cost
claim rather than an accuracy claim: machine-learned mitigation matches the
accuracy of conventional mitigation while substantially reducing its run-time
overhead, because a trained model needs only the unmitigated circuit at
inference time whereas extrapolation needs one execution per scale factor.
They also propose using learned models to imitate conventional mitigators as a
route to scalability.

Taken together, the literature on learned mitigation establishes three things
convincingly. Learned regressors can reduce residual error relative to
analytic extrapolation on the circuit families they are trained on
\cite{czarnik2021clifford,lowe2021unified,liao2024machine}. Training labels can
be obtained without full classical simulation, by near-Clifford construction,
echo evolution, or augmentation
\cite{czarnik2021clifford,babukhin2024echo,liao2025noiseagnostic}. And the
inference-time circuit cost of a trained mitigator can be lower than that of
extrapolation \cite{liao2024machine}.

What the literature does not establish is equally clear once the studies are
read side by side. The baselines differ: some papers compare mainly against the
unmitigated value \cite{bennewitz2022neural,zhukov2024quantum}, others against
zero-noise extrapolation with one particular choice of scaling and fit
\cite{lowe2021unified,liao2024machine}, and few compare against the full set of
analytic extrapolants --- linear, Richardson and exponential --- that a
practitioner actually chooses between. Cost accounting is usually informal:
Liao et al.\ are explicit that reduced overhead is the point
\cite{liao2024machine}, but most studies report accuracy at unstated total
circuit cost, and none of the studies reviewed here reports an
error-versus-circuit-evaluation curve from which a reader could read off which
method wins at a fixed budget. And with the exception of the extrapolation
experiments in \cite{liao2024machine} and the cross-device transfer test in
\cite{zlokapa2020deep}, the evaluation protocol is in-distribution: models are
tested on circuits drawn from the same family, depth range and noise profile as
their training data. Since the practical appeal of a learned mitigator is
precisely that it can be trained once and reused, this is the gap that matters
most.

\section{Quantum machine learning primitives and their known failure modes}\label{sec2:qml}

If a learned mitigator is to be built from a quantum model rather than a
classical one, the relevant primitives are variational circuits and quantum
kernels. Both are reviewed extensively elsewhere \cite{cerezo2021vqa}; what
follows concentrates on the properties that determine whether such a model can
be expected to help on a regression task of the kind studied in this thesis,
and on the substantial body of negative results that constrains the answer.

\subsection{Variational circuits, encodings and data reuploading}

A parameterised quantum circuit model computes
$f(\mathbf{x};\boldsymbol{\theta}) = \langle 0\rvert
U^{\dagger}(\mathbf{x},\boldsymbol{\theta})\, O \,
U(\mathbf{x},\boldsymbol{\theta}) \lvert 0\rangle$, where the classical input
$\mathbf{x}$ enters through gate rotation angles. P\'erez-Salinas et al.\
showed that interleaving encoding layers with trainable layers --- data
reuploading --- makes a single-qubit circuit a universal approximator, and that
the construction is the quantum analogue of a feature-map expansion repeated
at successive depths \cite{perezsalinas2020reuploading}. Schuld, Sweke and
Meyer supplied the structural explanation that matters most for this thesis: a
circuit whose encoding consists of Pauli rotations $e^{-i x_{j} P/2}$ computes a
truncated Fourier series in the data,
\begin{equation}
  f(x;\boldsymbol{\theta}) \;=\; \sum_{\omega \in \Omega} c_{\omega}(\boldsymbol{\theta})\, e^{i \omega x},
  \label{eq2:fourier}
\end{equation}
where the accessible frequency spectrum $\Omega$ is fixed by the eigenvalues of
the encoding generators and by the number of reuploading layers, while the
coefficients $c_{\omega}$ are what training adjusts
\cite{schuld2021effect}. Equation~\eqref{eq2:fourier} makes two predictions
that this thesis tests directly. Within the range of data on which the
frequencies were fitted, richer spectra permit closer approximation; outside
it, the model is not merely inaccurate but \emph{periodic}, so a
data-reuploading regressor extrapolates by repeating its learned pattern rather
than by continuing any trend. Angle-encoded quantum models therefore carry a
structural disadvantage on extrapolation tasks that no amount of training data
inside the original range can remove. Eq.~\eqref{eq2:fourier} therefore
suggests a specific mechanism for out-of-distribution failure of the
data-reuploading arm; Section~\ref{sec:ch5_bounded} tests it and finds that
periodicity accounts for the size of that failure but not for the failure itself,
which Section~\ref{sec:ch5_robustness} traces to inputs outside the training
range.

\subsection{Quantum kernels and the fidelity kernel}

The kernel view of these models is cleaner. Encoding $\mathbf{x}$ into a state
$\lvert \phi(\mathbf{x})\rangle = U_{\phi}(\mathbf{x})\lvert 0\rangle$ defines
the fidelity (or quantum embedding) kernel
\begin{equation}
  k(\mathbf{x},\mathbf{x}') \;=\; \bigl\lvert \langle \phi(\mathbf{x}) \vert \phi(\mathbf{x}')\rangle \bigr\rvert^{2}
  \;=\; \bigl\lvert \langle 0 \rvert U_{\phi}^{\dagger}(\mathbf{x})\,U_{\phi}(\mathbf{x}')\lvert 0\rangle \bigr\rvert^{2},
  \label{eq2:fidkernel}
\end{equation}
which is estimated on hardware either by the inversion test implied by the
right-hand side of Eq.~\eqref{eq2:fidkernel} or by a swap test. Havl\'i\v{c}ek
et al.\ introduced this construction with the IQP-style encoding that remains
the standard choice, and demonstrated both a variational classifier and a
kernel estimator on superconducting hardware \cite{havlicek2019supervised}.
Schuld later established the conceptual point that supervised models built from
parameterised circuits are kernel methods: training a variational circuit
explores a restricted region of a reproducing kernel Hilbert space whose kernel
is fixed by the encoding, so the encoding, not the ansatz, is where the
inductive bias lives \cite{schuld2021kernel}. Once a kernel matrix
$K_{ij} = k(\mathbf{x}_{i},\mathbf{x}_{j})$ has been estimated, regression is
classical kernel ridge regression,
\begin{equation}
  \widehat{y}(\mathbf{x}) \;=\; \sum_{i=1}^{N} \alpha_{i}\, k(\mathbf{x},\mathbf{x}_{i}),
  \qquad
  \boldsymbol{\alpha} \;=\; \bigl(K + \mu \mathbb{I}\bigr)^{-1}\mathbf{y},
  \label{eq2:krr}
\end{equation}
with $\mu$ a ridge parameter. This is the estimator used for the quantum arm
of the present benchmark, and its appeal is that the quantum device is
responsible only for the Gram matrix, so training is convex and there are no
barren plateaus to navigate.

Huang et al.\ analysed when such models can offer an advantage, introducing a
geometric quantity that compares the quantum and classical Gram matrices and
proving that a large separation is necessary for a prediction advantage; they
also showed how classical machine learning augmented by data can match quantum
models in many practically relevant regimes \cite{huang2021power}. Liu,
Arunachalam and Temme exhibited a learning problem, based on the discrete
logarithm, with a rigorous and noise-robust quantum speed-up for a
fidelity-kernel classifier \cite{liu2021rigorous}, establishing that advantage
is possible in principle while leaving open how far such a construction is
from anything a practitioner would encounter.

\subsection{Concentration, bandwidth, and the collapse of expressive kernels}

The most consequential practical result about fidelity kernels is that
expressive encodings destroy them. Thanasilp et al.\ proved that for broad
classes of encodings --- deep, expressive circuits, and also circuits subject
to global depolarising noise or measured with too many qubits traced out ---
the kernel entries concentrate exponentially in the number of qubits around a
fixed value,
\begin{equation}
  k(\mathbf{x},\mathbf{x}') \;=\; \bar{k} \;+\; \varepsilon(\mathbf{x},\mathbf{x}'),
  \qquad
  \mathrm{Var}\bigl[\varepsilon\bigr] \in \mathcal{O}\!\left(2^{-n}\right),
  \label{eq2:concentration}
\end{equation}
so that the Gram matrix approaches a constant matrix and, after
regularisation, the trained model approaches a trivial predictor unless the
number of shots grows exponentially \cite{thanasilp2024exponential}. This is
the kernel counterpart of the barren plateau phenomenon and it is not a
small-sample artefact: it is a property of the geometry of the embedding.

The practical remedy is bandwidth. Shaydulin and Wild observed that rescaling
the data before encoding, $\mathbf{x} \mapsto c\,\mathbf{x}$ with $c < 1$,
changes the classification accuracy of quantum kernel models dramatically and
can be the difference between near-random and competitive performance
\cite{shaydulin2022importance}. Canatar, Peters, Pehlevan, Wild and Shaydulin
then explained the effect within kernel learning theory: bandwidth controls the
spectral decay of the kernel's eigenvalues and therefore the
inductive bias of the model, and only at suitable bandwidth does generalisation
improve with qubit count instead of degrading \cite{canatar2023bandwidth}.
The consequence for experimental design is sharp: the bandwidth is a
hyperparameter that must be tuned, and a quantum kernel evaluated at the
default scaling $c = 1$ is not a fair representative of the method. It is also
a double-edged result, as Fl\'orez-Ablan et al.\ showed numerically: at the
bandwidths that optimise generalisation, fidelity and projected quantum kernels
closely resemble classical radial basis function kernels, and in the
small-bandwidth limit they reduce to low-order polynomial kernels ---
so the tuning that makes a quantum kernel work also makes it
classically imitable \cite{florezablan2025similarity}. Any claim of quantum
advantage from a bandwidth-tuned fidelity kernel must therefore be checked
against a classical kernel baseline.

\subsection{Trainability and generalisation}

For the variational arm, the dominant obstruction is the barren plateau:
McClean et al.\ showed that for sufficiently random parameterised circuits the
gradient of the cost function vanishes exponentially in the number of qubits,
so that optimisation becomes infeasible without structured ans\"atze or
informed initialisation \cite{mcclean2018barren}. The generalisation side is
more encouraging. Caro et al.\ proved bounds showing that the generalisation
error of a quantum model scales roughly as $\sqrt{T/N}$ in the number of
trainable gates $T$ and training samples $N$, so that good generalisation from
few samples is possible provided the model is not over-parameterised
\cite{caro2022generalization}. Crucially for this thesis, the same group
studied out-of-distribution generalisation explicitly for the task of learning
quantum dynamics and showed that models trained on product states can
generalise to entangled inputs for that structured problem
\cite{caro2023ood} --- which establishes that out-of-distribution behaviour is
a well-posed and analysable question for quantum models, not an unfair test.

Finally, the benchmarking literature itself has become sceptical. Schuld and
Killoran argued that framing quantum machine learning around advantage claims
has distorted the field's evaluation practices, and urged attention to
inductive bias and to careful comparisons instead \cite{schuld2022isquantum}.
Bowles, Ahmed and Schuld ran a large systematic benchmark of quantum models
against standard classical baselines and found that reported advantages
frequently fail to survive properly tuned classical comparison, and that
methodological choices --- hyperparameter search, data scaling, dataset
selection --- dominate the observed differences \cite{bowles2024better}. Their
recommended practice, which this thesis follows, is to fix the task, tune
every model including the classical ones, and report the comparison at matched
cost rather than at matched convenience.

\section{Quantum error correction and the decoding problem}\label{sec2:qec}

Mitigation buys accuracy on shallow circuits; error correction buys
scalability. The second track of this thesis concerns decoding, and the
literature there is instructive because it is a machine-learning subfield with
mature benchmarking conventions --- a fixed code, a fixed noise model, a single
reported figure of merit, and a universally accepted classical baseline.

\subsection{Surface codes and the decoding problem}

Kitaev's toric code introduced the idea of protecting quantum information in
the topology of a lattice of qubits with local stabiliser measurements
\cite{kitaev2003fault}, and Dennis, Kitaev, Landahl and Preskill developed the
planar variant into a full fault-tolerance framework, establishing the
mapping between decoding and a statistical-mechanical model and deriving
threshold estimates \cite{dennis2002topological}. Fowler et al.\ turned this
into an engineering blueprint, specifying the measurement circuits, the
lattice surgery operations and the resource estimates that made the surface
code the default architecture for superconducting processors
\cite{fowler2012surface}.

Operationally, a distance-$d$ surface code is measured repeatedly; each round
of stabiliser measurements yields a syndrome, and the differences between
consecutive rounds define \emph{detection events}. The decoder receives the
detection-event history $\mathbf{s}$ and must output a correction. The
statistically optimal rule is maximum-likelihood decoding over error classes,
\begin{equation}
  \widehat{c} \;=\; \arg\max_{c \,\in\, \{\text{logical classes}\}}\;
  \sum_{E \,:\, \sigma(E)=\mathbf{s},\; [E]=c} \Pr(E),
  \label{eq2:mld}
\end{equation}
where $\sigma(E)$ is the syndrome of error $E$ and $[E]$ its logical class.
Equation~\eqref{eq2:mld} is intractable in general, so practical decoding is a
search for accurate approximations that run faster than the syndrome
measurement cycle. Because decoding research depends on generating enormous
volumes of syndrome data under circuit-level noise, the field's progress is
tied to fast stabiliser simulation; Gidney's \textsc{Stim} made
million-shot circuit-level sampling routine and is now the standard data
generator \cite{gidney2021stim}.

\subsection{Matching decoders and their fast descendants}

The dominant classical approach maps decoding to minimum-weight perfect
matching (MWPM) on a graph whose vertices are detection events and whose edge
weights are negative log-likelihoods of the corresponding error mechanisms; the
matching itself is solved by Edmonds' blossom algorithm
\cite{edmonds1965paths}. This reduction, already present in
\cite{dennis2002topological}, is near-optimal for independent bit-flip and
phase-flip noise but suboptimal when errors are correlated between the two
lattices --- notably for $Y$ errors and for measurement errors with realistic
circuit-level structure. Higgott's \textsc{PyMatching} provided the open
implementation that made MWPM the universal baseline
\cite{higgott2022pymatching}, and the throughput problem was addressed by
Delfosse and Nickerson's union-find decoder, which achieves almost-linear
time in the number of syndrome bits at a modest cost in threshold by growing
and merging clusters instead of solving an exact matching
\cite{delfosse2021almost}. The practical decoder design space is thus a
trade-off surface between accuracy and latency, and any new decoder --- learned
or not --- is expected to report where it sits on that surface.

\subsection{Neural network decoders}

Machine-learned decoding was proposed to recover the accuracy that matching
gives up. Torlai and Melko trained a restricted Boltzmann machine as a
generative model of the error distribution conditioned on the syndrome and
showed it could decode the toric code near threshold \cite{torlai2017neural};
Varsamopoulos, Criger and Bertels trained feed-forward networks on small
surface codes, framing decoding as classification of the logical correction
\cite{varsamopoulos2018decoding}; and Baireuther et al.\ introduced recurrent
networks so that a decoder could consume an arbitrary number of measurement
rounds and, importantly, learn correlated error behaviour directly from data
rather than from a specified error model \cite{baireuther2018machine}. The
unifying argument in this literature is that a learned decoder needs no prior
estimate of physical error rates and can absorb error correlations that a
matching graph cannot express.

Recent work has made that argument quantitative. Varbanov et al.\ evaluated a
recurrent neural-network decoder on both simulated and experimental transmon
data for small-distance surface codes and found roughly $25\%$ lower logical
error rate than MWPM on the experimental data set, approaching
maximum-likelihood performance, with the gain traced specifically to better
handling of correlated defects such as those from $Y$ errors; they also
incorporated the soft (analog) readout information that a matching decoder
discards \cite{varbanov2025neural}. Lange et al.\ formulated decoding as a
graph classification problem and showed that a graph neural network trained
only on simulated experimental data can outperform a matching decoder under
circuit-level noise even when the matching decoder is given full knowledge of
the underlying error model \cite{lange2025data}. The current state of the art
is the transformer-based \textsc{AlphaQubit} decoder of Bausch et al., which
was pre-trained on simulated data, fine-tuned on experimental data, and
reported lower logical error rates than both matching and correlated-matching
baselines on Google's distance-3 and distance-5 surface code experiments
\cite{bausch2024learning}.

The methodological contrast with the mitigation literature is stark and
deliberate here. Every decoder paper cited above reports a single agreed
metric --- logical error rate per round --- against a single agreed baseline
--- MWPM, usually via \textsc{PyMatching} --- on a single agreed noise model,
and recent papers additionally report latency or throughput because a decoder
that cannot keep up with the syndrome cycle is useless regardless of accuracy.
Generalisation is also tested as a matter of course: decoders are evaluated at
code distances and error rates they were not trained on, and on experimental
data after training on simulation
\cite{varbanov2025neural,lange2025data,bausch2024learning}.

There is a second transferable lesson. Learned decoders win where the
classical baseline throws information away --- correlated $Y$ errors, analog
readout confidence, circuit-level error structure that a matching graph cannot
represent \cite{varbanov2025neural,lange2025data} --- and they do not win by
being more flexible in general. That is a useful diagnostic to carry into
mitigation: a learned mitigator should be expected to help where the analytic
estimator discards structure, for example by using several observables jointly
or by exploiting circuit metadata, and should not be expected to help merely
because it has more parameters than a polynomial fit. It also suggests the
right null hypothesis for a mitigation benchmark. Since the mitigation map for
global depolarising noise is a one-parameter rescaling
(Eq.~\eqref{eq2:depol}), a flexible regressor has little structure left to
exploit once that rescaling is captured, and the appropriate expectation is
parity with a well-tuned analytic estimator rather than a large improvement.

\subsection{Experimental status}

The experimental milestones establish that the whole enterprise is no longer
hypothetical. Google's surface code experiments demonstrated first a
distance-3 to distance-5 comparison and then, decisively, operation below
threshold: a distance-7 logical qubit with a logical error rate suppressed by
roughly a factor of two per distance increment, together with real-time
decoding at the measurement cycle rate \cite{google2025below}. In a different
hardware modality, Bluvstein et al.\ demonstrated a reconfigurable
neutral-atom processor performing logical operations on up to 48 logical
qubits with transversal gates and mid-circuit measurement
\cite{bluvstein2024logical}. Both results make decoder quality and decoder
latency the binding constraints on near-term logical performance, which is
precisely why learned decoders are studied so carefully and why the standards
of comparison in that subfield are worth importing into the mitigation
literature.

\section{Research gap}\label{sec2:gap}

The review above supports a narrow and specific claim: the case for
machine-learning-based quantum error mitigation has been made with an
evaluation protocol too weak to support the conclusions that are drawn from it.
Three deficiencies are distinguishable, and each has a concrete counterpart in
the design of this thesis.

\subsection{Inconsistent baselines}

There is no agreed reference method for
learned mitigation. Clifford data regression was introduced against the
unmitigated value and analysed in isolation \cite{czarnik2021clifford};
variable-noise CDR was compared against both ZNE and CDR, which is stronger,
but against one extrapolation order and one noise-scaling method
\cite{lowe2021unified}; the large-scale hardware benchmark of Liao et al.\ uses
digital ZNE as its single reference \cite{liao2024machine}; the
neural-network studies of Bennewitz et al.\ and Zhukov and Pogosov report
improvement over the noisy expectation value
\cite{bennewitz2022neural,zhukov2024quantum}; and application-driven work
reports whether the downstream optimisation succeeded rather than an error
metric against alternatives \cite{sack2024large}. Because ZNE is not one
method but a family --- linear, Richardson, and exponential fits, over
pulse-stretched \cite{kandala2019extends}, folded
\cite{giurgicatiron2020digital} or identity-inserted \cite{he2020zero} noise
scaling, each with its own bias--variance profile
\cite{endo2018practical,cai2021multiexponential} --- reporting a win against
``ZNE'' without specifying which member of the family was used, and without
tuning it, leaves the comparison unfalsifiable. The contrast with the decoding
literature, where every paper reports against MWPM under
\textsc{PyMatching} \cite{higgott2022pymatching} on a stated circuit-level
noise model, is the clearest evidence that the mitigation field's looseness is
a choice rather than a necessity. This thesis therefore fixes a single task and
evaluates nine methods on it --- the unmitigated baseline, linear, Richardson
and exponential extrapolation, three classical regressors, a fidelity-kernel
model and a data-reuploading circuit --- with identical data, identical shot
budgets and tuned hyperparameters for every arm, including the classical ones,
as Bowles et al.\ recommend \cite{bowles2024better}.

\subsection{Absent cost normalisation}

The theory of mitigation is a theory of
cost: unbiasedness is purchased with sampling overhead, and the fundamental
results are statements about how that overhead scales
\cite{takagi2022fundamental,quek2024exponentially}. Yet accuracy is routinely
reported without the corresponding budget. This matters because the methods
under comparison consume different numbers of circuit evaluations per
estimated observable. A three-point extrapolation consumes three; a classical
regressor trained on noise-scaled features consumes as many evaluations as it
has features; a fidelity-kernel or data-reuploading model consumes those
evaluations \emph{plus} the additional circuit executions needed to evaluate
its feature map or ansatz. An accuracy comparison that ignores this difference
credits the quantum models with resources it does not count. Liao et al.\ are
the exception that proves the rule: their central claim is explicitly about
run-time overhead \cite{liao2024machine}, and once cost is made the axis, the
conclusion about which method to prefer changes. This thesis therefore reports
error against circuit evaluations consumed, so that methods are ranked at
matched budget rather than at matched accuracy, and a model that ties on
accuracy while consuming more circuit executions is recorded as losing.

\subsection{No out-of-distribution testing}

The appeal of a learned mitigator
is amortisation --- train once, apply to many circuits. That appeal is only
justified if the learned map survives a change in the circuits it is applied
to, and almost no study tests this. Training circuits are constructed to
resemble the target \cite{czarnik2021clifford,lowe2021unified}, evaluation
draws from the training family, and the two available exceptions are partial:
the cross-device transfer test of Zlokapa and Gheorghiu, which found most of
the benefit did not transfer \cite{zlokapa2020deep}, and the extrapolation
experiments of Liao et al.\ \cite{liao2024machine}. Theory predicts that this
should be the decisive test for quantum models specifically. A
data-reuploading circuit is a truncated Fourier series in its inputs
\cite{schuld2021effect}, so its predictions are periodic beyond the training
range; a fidelity kernel concentrates unless its bandwidth is tuned
\cite{thanasilp2024exponential,shaydulin2022importance,canatar2023bandwidth},
and the tuning that makes it generalise also makes it resemble a classical
radial basis function kernel \cite{florezablan2025similarity}. That
out-of-distribution generalisation is nonetheless a tractable and meaningful
criterion for quantum models is established by Caro et al.\
\cite{caro2023ood}, alongside the sample-complexity bounds that govern
in-distribution behaviour \cite{caro2022generalization}. This thesis therefore
trains every model on circuits up to a fixed depth and tests all nine methods
on deeper circuits, reporting the degradation explicitly rather than the
in-distribution number alone.

Taken together, these three corrections define the contribution of this work.
It is not a new mitigation method, and it does not claim a quantum advantage;
it is a controlled, cost-normalised, distribution-shift-aware benchmark of
hybrid quantum-classical and purely classical mitigators on a common
noisy-to-ideal regression task, complemented by a machine-learned decoding
study that borrows its evaluation conventions from the more disciplined
decoding literature
\cite{varbanov2025neural,lange2025data,bausch2024learning}. The methodology
that implements these three corrections is set out in the next chapter.

\end{document}
